% Part: first-order-logic 
% Chapter: axiomatic-deduction 
% Section: provability

% verification of properties of provability needed for maximally 
% consistent sets in the completeness chapter.

\documentclass[../../../include/open-logic-section]{subfiles}

\begin{document}

\iftag{FOL}
      {\olfileid{fol}{axd}{prv}}
      {\olfileid{pl}{axd}{prv}}

\olsection{خواص \usetoken{S}{derivability}}

\begin{editorial}
نفهم في هذا القسم شرط الجدة في~\QR{} بالنسبة إلى فرضيات البرهان المستعملة
فعلًا، لا إلى أعضاء مجموعة محيطة لم تُستعمل فيه. وهذا اصطلاح الدعم المنتهي
الذي تعتمد عليه الرتابة والتراص.
\end{editorial}

\begin{prop}[الرتابة] إذا كان~$\Gamma \subseteq \Delta$ و
$\Gamma \Proves !A$، كان أيضًا~$\Delta \Proves !A$.
\end{prop}

\begin{proof}
كل~$\Gamma_0 \subseteq \Gamma$ منتهية هي أيضًا مجموعة جزئية منتهية
من~$\Delta$، ولذلك فإن~!!a{derivation} لـ$\Gamma_0 \Proves !A$ يبيّن
أيضًا~$\Delta \Proves !A$.
\end{proof}

\begin{prop}\ollabel{prop:provability}
$\Gamma \Proves !A$ إذا وفقط إذا كانت~$\Gamma\cup \{\lnot!A \}$ غير
متسقة.
\end{prop}

\begin{proof}
افترض أولًا~$\Gamma \Proves !A$. عندئذ توجد مجموعة منتهية
$\Gamma_0 \subseteq \Gamma$ بحيث~$\Gamma_0 \Proves !A$، ومن ثم
$\Gamma_0 \cup \{\lnot !A\} \Proves \lfalse$. ولذلك تكون
$\Gamma \cup \{\lnot !A\}$ غير متسقة. وبالعكس، إذا كانت
$\Gamma \cup \{\lnot !A\}$ غير متسقة، فثمة مجموعة منتهية
$\Gamma_0 \subseteq \Gamma$ بحيث
$\Gamma_0 \cup \{\lnot !A\} \Proves \lfalse$. وتعطي مبرهنة
الاستنباط~$\Gamma_0 \Proves \lnot !A \lif \lfalse$، وتعطي
\textbf{Ax0} الصيغة
$\Gamma_0 \Proves (\lnot !A \lif \lfalse) \lif !A$. وبـMP نحصل على
$\Gamma_0 \Proves !A$، ومن ثم~$\Gamma \Proves !A$ بالرتابة.
\end{proof}

\begin{prop}\ollabel{prop:provability-properties}
\begin{enumerate}
\item \ollabel{prop:provability-contr} إذا كان~$\Gamma \Proves !A$ و
$\Gamma \cup \{ !A\} \Proves \lfalse$، كانت~$\Gamma$ غير متسقة.

\item \ollabel{prop:provability-lnot} إذا كان
$\Gamma \cup \{!A\} \Proves \lfalse$، كان
$\Gamma \Proves \lnot !A$.

\item \ollabel{prop:provability-exhaustive} إذا كان
$\Gamma \cup \{!A\} \Proves \lfalse$ و
$\Gamma \cup \{\lnot !A\} \Proves \lfalse$، كان
$\Gamma \Proves \lfalse$.

\tagitem{prvOr}{\ollabel{prop:provability-lor-left} إذا كان
$\Gamma \cup \{!A\} \Proves \lfalse$ و
$\Gamma \cup \{!B\} \Proves \lfalse$، كان
$\Gamma \cup \{!A \lor !B\} \Proves \lfalse$.}{}

\tagitem{prvOr}{\ollabel{prop:provability-lor-right} إذا كان
$\Gamma \Proves !A$ أو~$\Gamma \Proves !B$، كان
$\Gamma \Proves !A \lor !B$.}{}

\tagitem{prvAnd}{\ollabel{prop:provability-land-left} إذا كان
$\Gamma \Proves !A \land !B$، كان~$\Gamma \Proves !A$ و
$\Gamma \Proves !B$.}{}

\tagitem{prvAnd}{\ollabel{prop:provability-land-right} إذا كان
$\Gamma \Proves !A$ و$\Gamma \Proves !B$، كان
$\Gamma \Proves !A \land !B$.}{}

\tagitem{prvIf}{\ollabel{prop:provability-mp} إذا كان
$\Gamma \Proves !A$ و$\Gamma \Proves !A \lif !B$، كان
$\Gamma \Proves !B$.}{}

\tagitem{prvIf}{\ollabel{prop:provability-lif} إذا كان
$\Gamma \Proves \lnot !A$ أو~$\Gamma \Proves !B$، كان
$\Gamma \Proves !A \lif !B$.}{}
\end{enumerate}
\end{prop}

\begin{proof}
\begin{enumerate}
\item اختر مجموعتين منتهيتين~$\Gamma_0,\Gamma_1 \subseteq \Gamma$
بحيث~$\Gamma_0 \cup \{!A\} \Proves \lfalse$ و
$\Gamma_1 \Proves !A$.

\begin{derivation}
1. & $\Gamma_0 \cup \{!A\} \Proves \lfalse$ & فرض \\
2. & $\Gamma_1 \Proves !A$ & فرض \\
3. & $\Gamma_0 \cup \Gamma_1 \Proves \lfalse$ & القطع \\
\end{derivation}

ولأن~$\Gamma_0 \subseteq \Gamma$ و$\Gamma_1 \subseteq \Gamma$، يكون
$\Gamma_0 \cup \Gamma_1 \subseteq \Gamma$، ومن ثم
$\Gamma \Proves \lfalse$.

\item اختر مجموعة منتهية~$\Gamma_0 \subseteq \Gamma$ بحيث
$\Gamma_0 \cup\{!A\} \Proves \lfalse$.

\begin{derivation}
1. & $\Gamma_0 \cup\{!A\} \Proves \lfalse$ & فرض \\
2. & $\Gamma_0 \Proves !A \lif \lfalse$ & مبرهنة الاستنباط \\
3. & $\Gamma_0 \Proves (!A \lif \lfalse) \lif \lnot !A$ & Ax0 \\
4. & $\Gamma_0 \Proves \lnot !A$ & MP 2, 3 \\
\end{derivation}

\item اختر مجموعتين منتهيتين~$\Gamma_0,\Gamma_1 \subseteq \Gamma$
بحيث~$\Gamma_0 \cup \{ !A \} \Proves \lfalse$ و
$\Gamma_1 \cup \{\lnot !A \} \Proves \lfalse$.

\begin{derivation}
1. & $\Gamma_0 \cup\{!A\} \Proves \lfalse$ & فرض \\
2. & $\Gamma_0 \Proves !A \lif \lfalse$ & مبرهنة الاستنباط \\
3. & $\Gamma_1 \cup \{\lnot !A \} \Proves \lfalse$ & فرض \\
4. & $\Gamma_0 \Proves (!A \lif \lfalse) \lif \lnot !A$ & Ax0 \\
5. & $\Gamma_0 \Proves \lnot !A$ & MP 2, 4 \\
6. & $\Gamma_0 \cup \Gamma_1 \Proves \lfalse$ & القطع 3، 5 \\
\end{derivation}

ولأن~$\Gamma_0 \subseteq \Gamma$ و$\Gamma_1 \subseteq \Gamma$، يكون
$\Gamma_0 \cup \Gamma_1 \subseteq \Gamma$. ومن ثم
$\Gamma \Proves \lfalse$.

% prop:provability-lor-left 
\tagitem{defOr}{}{
\iftag{probOr}{تمرين.}{تمرين.}}

% prop:provability-lor-right 
\tagitem{defOr}{}{\iftag{probOr}{تمرين.}{
افترض~$\Gamma_0 \Proves !A$.

\begin{derivation}
1. & $\Gamma_0 \Proves !A$ & فرض \\
2. & $\Gamma_0 \Proves !A \lif (!A \lor !B)$ & Ax0 \\
3. & $\Gamma_0 \Proves !A \lor !B$ & MP 1, 2 \\
\end{derivation}

إذن~$\Gamma \Proves !A \lor !B$. والبرهان في حالة
$\Gamma \Proves !B$ مماثل.}}

% prop:provability-land-left 
\tagitem{defAnd}{}{\iftag{probAnd}{تمرين}{
افترض~$\Gamma_0 \Proves !A \land !B$.

\begin{derivation}
1. & $\Gamma_0 \Proves !A \land !B$ & فرض \\
2. & $\Gamma_0 \Proves (!A\land !B) \lif !A $ & Ax0 \\
3. & $\Gamma_0 \Proves !A$ & MP 1, 2 \\
\end{derivation}

ومن ثم~$\Gamma \Proves !A$. ويبيّن~!!{derivation} مماثل أن
$\Gamma \Proves !B$.}}

% prop:provability-land-right
\tagitem{defAnd}{}{\iftag{probAnd}{تمرين.}{تمرين.}}

% prop:provability-mp 
\tagitem{defIf}{}{\iftag{probIf}{تمرين.}{نفترض
$\Gamma_0 \Proves !A$ و$\Gamma_0 \Proves !A \lif !B$ معًا.

\begin{derivation}
1. & $\Gamma_0 \Proves !A$ & فرض \\
2. & $\Gamma_0 \Proves !A \lif !B $ & فرض \\
3. & $\Gamma_0 \Proves !B$ & MP 1, 2 \\
\end{derivation}

ولأن~$\Gamma_0 \subseteq \Gamma$، يعني ذلك~$\Gamma \Proves !B$.}}

% prop:provability-lif 
\tagitem{defIf}{}{\iftag{probIf}{تمرين.}{افترض أولًا
$\Gamma \Proves \lnot !A$.

\begin{derivation}
1. & $\Gamma_0 \Proves \lnot !A$ & فرض \\
2. & $\Gamma_0 \Proves \lnot !A \lif (!A \lif !B) $ & Ax0 \\
3. & $\Gamma_0 \Proves !A \lif !B$ & MP 1, 2 \\
\end{derivation}

والبرهان في حالة~$\Gamma \Proves !B$ مماثل:

\begin{derivation}
1. & $\Gamma_0 \Proves !B$ & فرض \\
2. & $\Gamma_0 \Proves !B \lif (!A \lif !B) $ & Ax0 \\
3. & $\Gamma_0 \Proves !A \lif !B$ & MP 1, 2 \\
\end{derivation}}}

\end{enumerate}
\end{proof}

\begin{probtag}{probOr,probAnd,probIf}
أكمل برهان~\olref[fol][axd][prv]{prop:provability-properties}.
\end{probtag}

\begin{thm}[التعميم]\ollabel{thm:Generalization}
إذا كان~$\Gamma \Proves !A$ ولم يكن~$x$ حرًا في أي~!!{formula} في
$\Gamma$، كان~$\Gamma \Proves \lforall[x][!A]$.
\end{thm}

\begin{proof}
خذ برهانًا منتهيًا لـ~$!A$ من~$\Gamma$، ولتكن~$\Gamma_0$ مجموعة
!!{formula}s~$\Gamma$ التي استُعملت فيه. اختر~!!a{constant}~$c\in\mathcal E$ لا يقع
في أي صيغة من صيغ البرهان، ولا يساوي أي ثابت مميز استُعمل لتسويغ تطبيق
لـ~\QR{} فيه، وأحلّ~$c$ محل كل وقوع حر لـ~$x$ في كل خطوة. ولأن
$x$ غير حر في~$\Gamma_0$ تبقى خطوات الفرض كما هي. ويبيّن الاستقراء في
خطوات البرهان أن المتتالية الجديدة برهان لـ
$\Subst{!A}{c}{x}$ من~$\Gamma_0$: فالتعويض يحفظ أمثلة مسلمات الروابط
والمكمّمات ومسلّمات الهوية ذات الحدود المغلقة، ويحفظ MP وصورتي~\QR{}،
كما يحفظ شرط الجدة فيهما لأن~$c$ اختير جديدًا. ومن ثم
$\Gamma_0 \Proves \Subst{!A}{c}{x}$.

والآن نشتق:
\begin{derivation}
1. & $\Gamma_0 \Proves \Subst{!A}{c}{x}$ & كما تقدم \\
2. & $\Gamma_0 \Proves \Subst{!A}{c}{x} \lif
  (\ltrue \lif \Subst{!A}{c}{x})$ & \olref[prp]{ax:lif1} \\
3. & $\Gamma_0 \Proves \ltrue \lif \Subst{!A}{c}{x}$ & MP 1، 2 \\
4. & $\Gamma_0 \Proves \ltrue \lif \lforall[x][!A]$ & \QR{} 3 \\
5. & $\Gamma_0 \Proves \ltrue$ & \olref[prp]{ax:ltrue} \\
6. & $\Gamma_0 \Proves \lforall[x][!A]$ & MP 4، 5
\end{derivation}
وينطبق~\QR{} في الخطوة~4 لأن~$c$ لا يقع في~$\Gamma_0$ ولا في~$\ltrue$
ولا في~$!A$.
وأخيرًا تعطينا الرتابة~$\Gamma \Proves \lforall[x][!A]$.
\end{proof}

\begin{thm}\ollabel{thm:WeakGen}
(\emph{التعميم الضعيف على الثوابت}) إذا كان~$\Gamma \Proves !A$ وكان
$c$ !!a{constant} لا يظهر في~$\Gamma$، وُجد متغير~$x$ لا يظهر في~$!A$
بحيث~$\Gamma \Proves \lforall[x][\Subst{!A}{x}{c}]$، ولا يتضمن
البرهان~$c$.
\end{thm}

\begin{proof}
% تصحيح برهاني (0643-I2): عولجت مسلّمتا المكمّمات وصورتا QR بتجديد الثوابت ثم تبديل ثابتين تقابليًا.
خذ برهانًا منتهيًا~$\Pi$ لـ~$!A$ من~$\Gamma$، ولتكن~$\Gamma_0$
دعمه المنتهي. طبّق \olref[qua]{lem:qr-freshening} على~$\Pi$ مع مجموعة
المنع~$\{c\}$؛ فنحصل على برهان~$\Pi^*$ خاتمته~$!A$ ودعمه محتوى
في~$\Gamma_0$، ولا يساوي~$c$ أيٌّ من ثوابته المميزة. ولأن~$\Pi^*$
منتهٍ، اختر~$d\in\mathcal E$ لا يظهر في أي صيغة فيه ولا يساوي أي ثابت
مميز على عقدة~\QR{} فيه. واختر متغيرًا~$x$ لا يظهر في~$!A$.

لتكن~$\rho$ إعادة التسمية المنتظمة التقابلية للثوابت التي تبدّل~$c$ و~$d$
وتثبت سائر الرموز؛ ونضع~$!A_d=\rho(!A)$، ونكتب هذا اختصارًا~$!A[d/c]$
لأن~$d$ لم يظهر في~$!A$. يؤدي تطبيق~$\rho$ في
كل سطر من~$\Pi^*$ إلى برهان لـ~$\rho(!A)$ من~$\Gamma_0$. فدعم هذا البرهان لا يتغير
لأن~$c$ لا يظهر في~$\Gamma$؛ ومسلمات الروابط محفوظة. وتبقى مسلّمتا
المكمّمات حالتين منهما لأن~$\rho$ ترسل كل حد مغلق~$t$ إلى الحد المغلق
$\rho(t)$، وكذلك تحفظ مسلمات الهوية ذات الحدود المغلقة. وتحفظ~MP.
أما تطبيق~\QR{} ذو الثابت المميز~$a$، فإن~$a\ne c,d$ بفضل التجديد؛
ومن ثم تحفظ~$\rho$ مقدمته
وخاتمته وشرط الجدة، في كلتا صورتي~\QR{}. وهكذا يخلو البرهان الناتج من~$c$.

ضع~$!A_x=\Subst{!A}{x}{c}$. ولأن~$x$ و~$d$ لا يظهران في~$!A$، يكون
إحلال~$d$ محل~$x$ في~$!A_x$ هو بعينه~$!A_d=\rho(!A)=!A[d/c]$.
ومن البرهان السابق نشتق:
\begin{derivation}
1. & $\Gamma_0 \Proves !A_d$ & كما تقدم \\
2. & $\Gamma_0 \Proves !A_d \lif (\ltrue \lif !A_d)$
   & \olref[prp]{ax:lif1} \\
3. & $\Gamma_0 \Proves \ltrue \lif !A_d$ & MP 1، 2 \\
4. & $\Gamma_0 \Proves \ltrue \lif \lforall[x][!A_x]$
   & \QR{} 3، بالثابت المميز~$d$ \\
5. & $\Gamma_0 \Proves \ltrue$ & \olref[prp]{ax:ltrue} \\
6. & $\Gamma_0 \Proves \lforall[x][!A_x]$ & MP 4، 5
\end{derivation}
ويصح تطبيق~\QR{} لأن~$d$ لا يظهر في~$\Gamma_0$ ولا في~$\ltrue$ ولا في
$!A_x$. وبالرتابة نحصل على
$\Gamma \Proves \lforall[x][\Subst{!A}{x}{c}]$ ببرهان لا يتضمن~$c$.
\end{proof}

\begin{explain}
هدفنا أن نستبدل باشتراط~\olref{thm:WeakGen} ألا يظهر~$x$ في~$!A$
إطلاقًا اشتراطين أضعف: ألا يكون~$x$ حرًا في~$!A$، وأن يكون~!!{free for}
$c$ فيها. ويمكن تحقيق ذلك بوضوح بتغيير~!!{variable} مقيد، ولذلك نبدأ
بإثبات هذا الأمر.
\end{explain}

\begin{lem}
\ollabel{lem:free-for}
إذا كان~$x$ حرًا لـ$c$ في~$!A$، ولم يكن~$y$ حرًا في~$!A$، كان~$x$
!!{free for}~$y$ في~$\Subst{!A}{y}{c}$.
\end{lem}

\begin{proof}
إذا لم يكن~$x$ !!{free for}~$y$ في~$\Subst{!A}{y}{c}$، وقع ظهور حر
ما لـ$y$ في~$\Subst{!A}{y}{c}$ داخل نطاق سور~$\lforall[x]$. لكن~$y$
ليس حرًا في~$!A$ بالفرض، ولذلك تأتي كل هذه الظهورات من التعويض
$\Subst{}{y}{c}$. ومن ثم يقع ظهور ما لـ$c$ داخل نطاق~$\lforall[x]$،
ولا يكون~$x$ !!{free for}~$c$ في~$!A$.
\end{proof}

\begin{lem}[تغيير المتغير المقيد]
\ollabel{lem:ChangeBdVar}
إذا لم يكن~$x$ ولا~$y$ حرًا في~$!A$، وكان كلاهما~!!{free for}~$c$ في
$!A$، كان
$\Proves \lforall[x][\Subst{!A}{x}{c}] \equiv
\lforall[y][\Subst{!A}{y}{c}]$، ولا يتضمن البرهان~$c$.
\end{lem}

\begin{proof}
بموجب الفرضيات، يكون~$y$ !!{free for}~$c$ في~$!A$، ولا يكون~$x$ حرًا
في~$!A$؛ ومن ثم، باللمّة السابقة~\olref{lem:free-for}، يكون~$y$
!!{free for}~$x$ في~$\Subst{!A}{x}{c}$. ولذلك تكون
$\lforall[x][\Subst{!A}{x}{c} \lif
\Subst{!A}{x}{c} \Subst{}{y}{x}]$ مسلّمة (\textbf{Ax1}). وبما أن~$x$
ليس حرًا في~$!A$، يكون
$\Subst{!A}{x}{c} \Subst{}{y}{x} = \Subst{!A}{y}{c}$، ومن ثم
\[
\Proves \lforall[x][\Subst{!A}{x}{c}] \lif \Subst{!A}{y}{c}.
\]
وتعطي مبرهنة الاستنباط
$\lforall[x][\Subst{!A}{x}{c}] \Proves \Subst{!A}{y}{c}$. ولأن~$y$
ليس حرًا في~$!A$، فهو أيضًا ليس حرًا في
$\lforall[x][\Subst{!A}{x}{c}]$؛ ومن ثم يعطينا التعميم، مع اختيار
ثابته المميز الجديد من~$\mathcal E\setminus\{c\}$،
$\lforall[x][\Subst{!A}{x}{c}] \Proves
\lforall[y][\Subst{!A}{y}{c}]$، وتعطي مبرهنة الاستنباط مرة أخرى
$\Proves \lforall[x][\Subst{!A}{x}{c}] \lif
\lforall[y][\Subst{!A}{y}{c}]$. وبرهان الاتجاه
$\Proves \lforall[y][\Subst{!A}{y}{c}] \lif
\lforall[x][\Subst{!A}{x}{c}]$ متناظر تمامًا، مع اختيار الثابت المميز
بالطريقة نفسها. وجميع الصيغ والخطوات السابقة خالية من~$c$، كما أن بناء
مبرهنة الاستنباط لا يستحدثه؛ فتتبع النتيجة بـTaut ببرهان يخلو من~$c$.
\end{proof}

\begin{thm}[التعميم القوي على الثوابت]
\ollabel{thm:StrongGen}
إذا كان~$\Gamma \Proves !A$، ولم يظهر~!!{constant}~$c$ في~$\Gamma$،
ولم يكن~$x$ حرًا في~$!A$ لكنه كان~!!{free for}~$c$ فيها، كان
$\Gamma \Proves \lforall[x][\Subst{!A}{x}{c}]$.
\end{thm}

\begin{proof}
بما أن~$\Gamma \Proves !A$ و$c$ لا يظهر في~$\Gamma$، يوجد بالتعميم
الضعيف~!!a{variable}~$y$ لا يظهر في~$!A$ بحيث
$\Gamma \Proves \lforall[y][\Subst{!A}{y}{c}]$. ولأن~$y$ لا يظهر في
$!A$ إطلاقًا، فهو ليس حرًا فيها، وهو أيضًا~!!{free for}~$c$ فيها. وإذا
كان~$x$، كما تفترض المبرهنة، غير حر في~$!A$ و!!{free for}~$c$ فيها،
تحققت شروط تغيير~!!{variable} مقيد، فيكون
$\Proves \lforall[x][\Subst{!A}{x}{c}] \equiv
\lforall[y][\Subst{!A}{y}{c}]$، ومنه
$\Gamma \Proves \lforall[x][\Subst{!A}{x}{c}]$.
\end{proof}

\begin{thm}
\ollabel{thm:provability-quantifiers}
\begin{tagenumerate}{prvEx,prvAll}
\tagitem{prvEx}{إذا كان~$\Gamma \Proves !A(t)$، كان
$\Gamma \Proves \lexists[x][!A(x)]$.}{}
\tagitem{prvAll}{إذا كان~$\Gamma \Proves \lforall[x][!A(x)]$، كان
$\Gamma \Proves !A(t)$.}{}
\end{tagenumerate}
\end{thm}

\begin{proof}
\begin{tagenumerate}{prvEx,prvAll}
%%Need axioms for exists
\tagitem{prvEx}{تمرين.}{}
\tagitem{prvAll}{افترض~$\Gamma_0 \Proves \lforall[x][!A(x)]$.

\begin{derivation}
1. & $\Gamma_0 \Proves \lforall[x][!A(x)]$ & فرض \\
2. & $\Gamma_0 \Proves \lforall[x][!A(x)] \lif \Subst{!A}{t}{x}$ & Ax1 \\
3. & $\Gamma_0 \Proves \Subst{!A}{t}{x}$ & MP 1, 2 \\
\end{derivation}}{}

\end{tagenumerate}
\end{proof}

\end{document}
